% Clase 23 --- La incoherencia de la ley de Ampère (§4.2).
% El docente dibuja el cable que carga un condensador y elige DOS superficies
% distintas con el mismo borde. La figura tiene que dejar ver que el borde es
% el mismo (la amperiana C) y que sin embargo una superficie es atravesada por
% la corriente y la otra no.
\begin{center}
\begin{tikzpicture}[scale=0.95, transform shape]

  % ---- superficie S2: la "bolsa" que se cuela entre las placas ----
  \fill[figamber!14]
    (0,1.35) .. controls (1.4,2.3) and (2.4,2.2) .. (2.1,0)
    .. controls (2.4,-2.2) and (1.4,-2.3) .. (0,-1.35)
    -- cycle;
  \draw[figamber, line width=0.9pt]
    (0,1.35) .. controls (1.4,2.3) and (2.4,2.2) .. (2.1,0)
    .. controls (2.4,-2.2) and (1.4,-2.3) .. (0,-1.35);

  % ---- cable y condensador ----
  \draw[figline] (-3.1,0) -- (1.6,0);
  \draw[figline] (2.6,0) -- (4.6,0);
  \draw[figline, line width=1.6pt] (1.6,-1.0) -- (1.6,1.0);
  \draw[figline, line width=1.6pt] (2.6,-1.0) -- (2.6,1.0);
  \node[figlbl, figred, anchor=east] at (1.5,0.72) {$+Q$};
  \node[figlbl, figblue, anchor=west] at (2.7,0.72) {$-Q$};

  \draw[figvecres] (-2.9,0.42) -- (-2.0,0.42);
  \node[figlbl, figred, anchor=south] at (-2.45,0.5) {$I$};
  \draw[figvecres] (3.5,0.42) -- (4.4,0.42);

  % ---- superficie S1: el disco plano ----
  \fill[figblue!14] (0,0) ellipse (0.3 and 1.35);

  % ---- amperiana C (borde común de las dos superficies) ----
  \draw[figline, line width=1.1pt] (0,0) ellipse (0.3 and 1.35);
  \node[figlbl, figblue, anchor=east] at (-0.42,-1.15) {$C$};

  % ---- rótulos ----
  \node[figlbl, figblue, anchor=east, align=right] at (-1.05,1.9)
    {$S_1$: corta el cable\\ $\Rightarrow \mu_0 I$};
  \draw[figvecaux, figblue] (-1.0,1.75) -- (-0.18,0.95);

  \node[figlbl, figamber, anchor=west, align=left] at (2.95,2.35)
    {$S_2$: pasa entre las placas\\ $\Rightarrow 0$};
  \draw[figvecaux, figamber] (2.9,1.95) -- (2.3,1.0);

  \node[figred, font=\Large\bfseries] at (1.15,-2.25) {?};
  \node[figlbl, figred, anchor=west] at (1.45,-2.25)
    {mismo borde, dos respuestas};

\end{tikzpicture}\\[0.5em]
{\small\itshape La ley de Ampère habla de \enquote{la corriente que atraviesa
una superficie cuyo borde es $C$}, pero no dice cuál: cualquiera debería
servir, y ahí se cae. Las dos superficies dibujadas tienen exactamente el mismo
borde ---la amperiana $C$, que abraza el cable--- y sin embargo el disco plano
$S_1$ es perforado por la corriente $I$, mientras que la bolsa $S_2$, que se
mete por el hueco entre las placas, no es atravesada por ninguna carga en
movimiento. El miembro izquierdo $\oint_C \mathbf{B}\cdot d\boldsymbol{\ell}$
es uno solo, así que el derecho no puede valer $\mu_0 I$ y $0$ a la vez. La
pista de dónde está la falla la da el condensador: la carga $Q$ de las placas
\textbf{crece con el tiempo}, o sea que el sistema no es estacionario, que es
justamente la hipótesis bajo la cual se dedujo la ley de Ampère. Lo que sí
ocurre en el hueco es que el campo eléctrico aumenta, y ése será el término que
Maxwell agregue para tapar el agujero.}
\end{center}
